% BHU PYQ -- Transcribed & Corrected
% Paper: Madhava Mathematics Competition 2024
% Organized by: Bhaskaracharya Pratishthana, Pune

\documentclass[11pt,a4paper]{article}
\usepackage[a4paper,top=2.5cm,bottom=2.5cm,left=2.8cm,right=2.8cm,headheight=14pt]{geometry}
\usepackage{amsmath,amssymb}
\usepackage[T1]{fontenc}
\usepackage[utf8]{inputenc}
\usepackage{lmodern,microtype,fancyhdr,enumitem,hyperref,lastpage,parskip}

\hypersetup{colorlinks=true,urlcolor=black,linkcolor=black,pdftitle={Madhava Mathematics Competition 2024}}
\pagestyle{fancy}\fancyhf{}
\fancyhead[L]{\small\textit{Madhava Mathematics Competition}}
\fancyhead[R]{\small\textit{UG Mathematics 2024}}
\fancyfoot[C]{\small Page \thepage\ of \pageref{LastPage}}

\newlist{parts}{enumerate}{2}
\setlist[parts,1]{label=\arabic*.,leftmargin=*,itemsep=6pt,topsep=2pt}

\begin{document}

\begin{center}
  {\large\bfseries MADHAVA MATHEMATICS COMPETITION 2024}\\[2pt]
  {\normalsize (A Mathematics Competition for Undergraduate Students)}\\[4pt]
  {\small Organized by: Bhaskaracharya Pratishthana, Pune}\\[6pt]
  \rule{0.8\linewidth}{0.5pt}\\[6pt]
  {\normalsize Time: 12.00 noon to 3.00 p.m. \hfill Maximum Marks: 100}
\end{center}

\rule{\linewidth}{0.4pt}\smallskip
\noindent\textit{N.B.: Part I carries 20 marks, Part II carries 30 marks, and Part III carries 50 marks.}
\bigskip

\noindent{\large\bfseries PART I (10 Questions $\times$ 2 Marks = 20 Marks)}
\begin{parts}
  \item Let $a < b < c$ be three positive integers in geometric progression whose common ratio is an integer. If the arithmetic mean of $a, b, c$ is $b + 2$, then the value of $a$ is: \\
  (A) 24 \quad (B) 12 \quad (C) 6 \quad (D) 4

  \item Triangles are formed by joining the vertices of a regular polygon of 12 sides. The number of triangles having no side in common with the polygon is: \\
  (A) 124 \quad (B) 220 \quad (C) 96 \quad (D) 112

  \item If $1, \alpha_1, \alpha_2, \dots, \alpha_{n-1}$ are the $n$-th ($n > 2$) roots of unity, then $\sum_{k=1}^{n-1} k \alpha_k$ equals: \\
  (A) $\frac{n}{2}$ \quad (B) $\frac{n(n-1)}{2}$ \quad (C) $-\frac{n}{1 - \alpha}$ \quad (D) $-\frac{n}{2}$

  \item If $f(x) = \begin{cases} ax + b, & x < -1 \\ 4x^2 + 1, & -1 \le x \le 1 \\ 2x - 1, & x > 1 \end{cases}$ is continuous everywhere, then which value can $ab$ never take? \\
  (A) $-3$ \quad (B) $-1$ \quad (C) $0$ \quad (D) $3$

  \item Let $f: [1, 3] \to \mathbb{R}$ be a continuous function. If $f(x)$ takes rational values for all $x$ and $f(2) = 10$, then $f(1.5) =$: \\
  (A) 7.5 \quad (B) 10 \quad (C) 1.5 \quad (D) 5

  \item If $f(x) = (x + 1)^2 - 1$ for $x \ge -1$, then the set $\{x : f(x) = f^{-1}(x)\}$ is: \\
  (A) $\{0\}$ \quad (B) $\{0, -1\}$ \quad (C) $\mathbb{R}$ \quad (D) $\emptyset$

  \item The number of real solutions of $1 + |x - 1| = e^x (x - 2)$ is: \\
  (A) 0 \quad (B) 1 \quad (C) 2 \quad (D) 3

  \item In how many ways can you choose an odd number of objects from $n$ objects? \\
  (A) $2^{n-1}$ \quad (B) $2^n$ \quad (C) $2^{n+1}$ \quad (D) $n$

  \item The family of circles with center $(a, 0)$ and radius $a$ satisfies the differential equation: \\
  (A) $2xy y' + y^2 - x^2 = 0$ \quad (B) $2xy y' + x^2 - y^2 = 0$ \quad (C) $2xy + x - y = 0$ \quad (D) $2xy' + x^2 + y = 0$

  \item The series $\sum_{n=1}^\infty \frac{n x \cos(n x) + \sin(n x) - e^{-|x|}}{n^3 + 1}$: \\
  (A) converges pointwise everywhere on $\mathbb{R}$, but not uniformly \quad (B) does not converge at any point \quad (C) converges uniformly on $\mathbb{R}$ \quad (D) converges only at 0
\end{parts}

\bigskip
\noindent{\large\bfseries PART II (5 Questions $\times$ 6 Marks = 30 Marks)}
\begin{parts}
  \item Let $f: [0, 1] \to \mathbb{R}$ be a continuous function with $\int_0^1 x(x-1)^2 f(x) \, dx = 0$. Prove that there exists $c \in (0, 1)$ such that $\int_0^c x^2 f(x) \, dx = c \int_0^c x f(x) \, dx$.
  \item Let $f(x) = x^2 + 4x + 3$ and $g(x) = 2x + 5$. Show that $\gcd(f(n), g(n)) = 1 \text{ or } 3$ as $n$ varies over integers.
  \item Let $a_1, a_2, \dots, a_n$ be real numbers. Prove that $\sum_{i,j=1}^n e^{a_i a_j} \ge 0$.
  \item Suppose for $n \ge 4$, $x_1, x_2, \dots, x_n$ are non-negative integers such that after removing any one of them, the remaining numbers can be partitioned into two sets with equal sums. (i) If $n$ is even, prove $x_i = 0$ for all $i$. (ii) Does the same hold when $n$ is odd?
  \item Give an example of a function $f: \mathbb{R} \to \mathbb{R}$ which is continuous at $a$ if and only if $a$ is a positive integer.
\end{parts}

\bigskip
\noindent{\large\bfseries PART III (4 Problems = 50 Marks)}
\begin{parts}
  \item Suppose $\alpha, \beta$ are roots of $x^2 + bx + c$. Let $p \ge 2$ be a prime such that $p \mid b, p \mid c$, but $p^2 \nmid c$.
  (i) Find largest power of $p$ dividing $\alpha^2 + \beta^2$. [2]
  (ii) Find largest power of $p$ dividing $\alpha^3 + \beta^3$. [3]
  (iii) For any $n$, find largest power of $p$ dividing $\alpha^{2^n} + \beta^{2^n}$. [8]

  \item (i) Show that $x^6 - 6x^5 + 1 = 0$ has a unique root in $(0, 1)$. [3]
  (ii) Show that every complex non-real root $z$ satisfies $|z| < z_0$, where $z_0$ is the root in $(0, 1)$. [10]

  \item (i) Suppose $A, B$ are $n \times n$ complex matrices with $AB^2 - B^2A = B$. Show $B$ is nilpotent. [9]
  (ii) Does there exist $A$ for a given nilpotent $B$ such that $AB^2 - B^2A = B$? [3]

  \item Let $A, B, C, D$ be vertices of a quadrilateral inscribed in the unit circle, and $P$ be inside $ABCD$. Show that $|PA||PB||PC||PD| \le \frac{27}{16}$. [12]
\end{parts}

\vfill\rule{\linewidth}{0.4pt}
\begin{center}\small\href{https://bhu-exam.vercel.app/}{Click here to visit our website}\end{center}
\end{document}
